\section{The outgoing one-sided symbols really have integral lifts}
\label{sec:outgoing-realization}

\status{This proves the additional outgoing realization used by the cloud
test. Its polynomial family is \texttt{monPh}, which needs a separate
determinant extraction before the earlier SY correction applies. The
actual numerical hypotheses with $C\ge100$ suffice.}

\paragraph{OR1--OR2: minimum weight forces a determinant.}
Work first over any field. For signed letters, let $v_S(P)$ be the least
$t$-exponent after a generic conjugate of
$\operatorname{diag}(1,\ldots,1,t)$ acts on the letters in $S$.
Vectors transform directly and covectors inversely. Choose $r$ distinct
vector letters $A\subset S$, put $D_A=\det(X_A)$, and let
$B$ be the total degree in the other vector letters of $S$.
For a letter-multihomogeneous polynomial,
\[
 v_S(P)\ge\theta>B
 \quad\Longrightarrow\quad D_A^{\theta-B}\mid P
 \qquad(\theta\in\mathbb Z).                              \tag{OR.1}
\]
Indeed, modulo the prime polynomial $D_A$, the generic selected matrix has
rank $r-1$. Choose a basis whose last direction it annihilates. The selected
vectors are then fixed by the cocharacter; the other vectors contribute
degree at most $B$ in $t$, and covectors only nonpositive powers.
The lower bound $\theta>B$ forces the specialized polynomial to vanish.
Thus $D_A\mid P$.

For any other subset $S'$, transforming only its selected rows is a
rank-one matrix update. Its determinant is affine in $t$, with minimum
exponent at most $1$. Consequently
\[
 v_{S'}(D_A^mH)\le v_{S'}(H)+m.                           \tag{OR.2}
\]
Iterating proves (OR.1), since division decreases only the selected
letters' degrees. Over an integral domain the quotient remains integral:
$D_A^m$ has unit leading coefficient, so polynomial division over the
domain gives the same zero remainder as over its fraction field.

\paragraph{OR3--OR5: extracting the actual segment divisor.}
Write $k=r+K$, $Q=p^{G'}$, $c=c_P$, $N=|TT|$, and
$\alpha_*=\max_{0\le j<J}\alpha_j$. The exact exponents are
\[
\begin{aligned}
m_1&=(2Q+3)c+kQ+c,&m_2&=c\alpha_*+2c,\\
m_*&=(N_1+1)\bigl((2Q+3)c+kQ+c\alpha_*+3c\bigr).
\end{aligned}                                                   \tag{OR.3}
\]
For each late segment letter $x$, let $A_x$ be its $r$ most recent even
selected predecessors, and let $D_x$ be their plane determinant.
Let $D_{{\rm old},j}$ be the preceding-frame determinant. Extract
\[
                 Dv_j=\prod_{x\ {\rm late}}D_x^{m_1}D_{{\rm old},j}^{m_2}.
                                                               \tag{OR.4}
\]
The total extraction cost is at most $|{\rm letters}_j|m_1+m_2\le m_*$.
Formula (OR.2) controls \emph{every} prefix after each extraction; the
other-vector degrees can only decrease. Thus a starting minimum
$B_x+m_*$ suffices for all tasks. The $D_{{\rm old},j}$ task uses the prefix before
the first late letter, containing exactly the preceding frame, so its
other-vector degree is zero.

The outgoing degree assumption is only $d_y\le1.01\,\operatorname{levDeg}(y)$
on the preceding frame and intermediate letters. An even predecessor
outside $A_x$ lies in these sets and is $2h$ positions before the oldest
member of $A_x$, for a distinct integer $h\ge1$. Hence, writing $B_o$ for
that oldest member's normalized degree,
\[
 B_x\le2.02B_o\sum_{h\ge1}p^{-2h}
       =\frac{2.02B_o}{p^2-1}\le\frac{B_o}{80}
       \qquad(p\ge13).                                    \tag{OR.5}
\]
There is no lower approximation assumption on $d_y$.

The actual hypotheses give $r,k\ge2$, $Y\ge r$, $N\ge1$, $c\ge1$,
$N_1\le k$, $\alpha_*\le1+rkQ$, and
$a\ge Ck^3r^2NQ^5cY$. The frame span gives $B_o\ge a/Q$. Therefore
\[
 m_*\le8rk^2Qc\le B_o/100,\qquad
 s=\lfloor B_o/(Yk^2)\rfloor\le B_o/8
 \quad(C\ge100).
\]
The prescribed minimum $B_o/2-s$ exceeds $B_x+m_*$.
Division thus leaves exactly the zero-gap plane/normal conditions
\texttt{MGen0}.

\paragraph{OR6: passing from the entire constrained space.}
An element of \texttt{XX'} with fixed whole-letter degrees is a finite
integral sum of the plane/normal products just considered. To see this,
project its defining product span to those degrees: the generic
substitutions preserve each letter degree, so every component retains all
minimum/maximum inequalities. The final order-support intersection can
be forgotten for this inclusion. After division, IP6 and IP11 give
\[
 \xi_{ij}=Dv_j\,\xi'_{ij},\qquad
 \phi_j(\xi'_{ij})\in\mathcal S,
\]
where $\mathcal S$ is the algebra of finite sums of key symbols with their
prescribed normal orders.

\paragraph{OR7: the exact extra multiplier.}
Use $D_s,\Delta_j,R_0$ for the already substituted top determinants,
frame determinants and $R_0=\prod_j\Delta_j^{\alpha_j}$.
Set
\[
 {\cal F}_c=R_0^c\prod_j\Delta_j^c\prod_{s\in TT}D_s^{(2Q+3)c},
 \qquad {\cal D}=\prod_{s\in TT}D_s^{kQ}.
\]
Entrywise Frobenius gives, inside the integral tangential algebra,
\[
 \phi_j(D_x)=D_{\bar x}^{p^{z_j-x}}+pE_x,\qquad
 \phi_j(D_{{\rm old},j})=\Delta_{j-1}^{p^{z_j-z_{j-1}}}+pE_j.
\]
The binomial theorem gives
$(D^P+pE)^m\in D^{m-c}\mathcal T+p^c\mathcal T$ when $P\ge1,m\ge c$.
The late intervals partition one full period, so each $s\in TT$ occurs
once. Combining the factors, and absorbing
$\prod_j\Delta_j^{c(\alpha_*-\alpha_j)}$, proves the literal
\texttt{RestPh} identity
\[
                    \prod_j\phi_j(Dv_j)={\cal F}_c{\cal D}H+p^cY,
 \qquad H\in\mathcal T,\quad Y\text{ integral}.             \tag{OR.6}
\]

\paragraph{OR8: the outgoing lift and its precision.}
For the actual fine label $X^\mu$, put
$G=\sum_i\prod_j\phi_j(\xi'_{ij})$. Then
\[
 \sigma=\operatorname{monPh}
       =(X^\mu{\cal D}GH){\cal F}_c+p^cX^\mu GY.            \tag{OR.7}
\]
Here $G\in\mathcal S$ is integral, and SY4 puts
$X^\mu{\cal D}$ in $\mathcal S$. Take the normal-degree-$o$ component.
SY5's finite correction applied to the first term gives an integral
$F_o$ such that
\[
 \Phi_gF_o\equiv w^o\sigma\pmod{w^{o+c}},\qquad
 \Phi_{g'}F_o\equiv0\pmod{w^{o+c}}
 \quad(\Pi(g')\ne\Pi(g)).                                 \tag{OR.8}
\]
The second term contributes $w^op^c=w^{o+2c}$, so it respects this
precision. Projection preserves any prescribed position multidegree.
These are all clauses of the actual \texttt{Realisable} conclusion.
The result closes the outgoing one-sided application under the supplied
Type-II hypotheses; it asserts no global supply of those hypotheses.
